\section{现实检验：一次真实部署的日志分析}
\label{sec:reality}

前述章节描述了 WhatyTerm 的设计蓝图。本节转向一个更关键的问题：\emph{它真的这样工作吗？} 我们对系统在作者机器上一次真实部署的完整历史记录进行分析。数据来源为系统自身持久化的 SQLite 历史表（\texttt{history}）与自主执行引擎的输出文件，均为运行时自动记录、非事后构造。

\subsection{数据集}

该部署跨越 2026 年 7 月 17 日至 20 日，共三天。历史表含 62 条记录，其中系统事件 9 条、终端输出 7 条、\textbf{自动决策（\texttt{ai\_decision}）46 条}。这些决策分布于四个会话，其目标均与“将本 harness 构建为一篇论文/软件工程任务”相关。表~\ref{tab:sessions} 给出各会话的决策计数。

\begin{table}[t]
  \centering\small
  \caption{四个会话的自动决策计数（共 46 次）。}
  \label{tab:sessions}
  \begin{tabular}{@{}lc@{}}
    \toprule
    会话（目标摘要） & 决策数 \\
    \midrule
    afad3bb2（自主开发迭代） & 36 \\
    83993863（论文写作） & 4 \\
    630f0238（arXiv 论文构建） & 3 \\
    8bf5eba6（论文写作） & 3 \\
    \midrule
    合计 & 46 \\
    \bottomrule
  \end{tabular}
\end{table}

\subsection{发现一：规则层 100\% 垄断决策}

最直接的发现是：\textbf{全部 46 次自动决策，无一例外地由规则层（标记为“预判断自动操作”）做出；LLM 兜底层的调用次数为零。} 按 \S\ref{sec:method-decision} 的定义，规则命中率
\[
\rho = \frac{N_{\text{pre}}}{N_{\text{pre}}+N_{\text{llm}}} = \frac{46}{46+0} = 100\%.
\]
在原始设计叙事里，高 $\rho$ 被视作“以零 Token 处理多数常见状态”的成本优势。但一个 $100\%$ 的 $\rho$ 意味着截然不同的事情：\emph{系统从未认为任何一次情形复杂到需要动用 LLM 的语义判断}。设计中作为“智能兜底”的第二级决策，在整个部署周期内是完全的死代码路径。

\subsection{发现二：87\% 的动作只是“继续”}

进一步看 46 次决策的动作构成（表~\ref{tab:actions}）：\textbf{“继续”占 40 次（87\%）}，重启 CLI（\texttt{claude -c} 或 \texttt{claude}）占 5 次，而携带真实语义的任务指令——形如“请继续完成第 7/15 个任务：实现……”——\textbf{仅出现 1 次（2\%）}。

\begin{table}[t]
  \centering\small
  \caption{46 次决策的动作构成。}
  \label{tab:actions}
  \begin{tabular}{@{}lcc@{}}
    \toprule
    动作 & 次数 & 占比 \\
    \midrule
    “继续” & 40 & 87\% \\
    重启 CLI（\texttt{claude -c}/\texttt{claude}） & 5 & 11\% \\
    真实任务指令 & 1 & 2\% \\
    \bottomrule
  \end{tabular}
\end{table}

这组数字刻画了一个令人不安的运行画像：在近九成的时间里，这个“自主开发编排系统”所做的，不过是反复向终端敲入“继续”二字。它既没有理解代理当前卡在何处，也没有判断这一声“继续”是否真的推动了任务——它只是在“看到空闲提示符”这一条件反射下，机械地重复同一个动作。

\subsection{发现三：自主引擎从未真正交付}

与交互式决策的空转相呼应，同期的三层自主引擎（RalphEngine）表现同样失败。我们检查其 headless 执行的输出文件：早期批次（十余个）\textbf{输出为空}，仅有一行完成标记；后期批次则直接写入 \texttt{Execution error}——这是被调用的 CLI 启动即失败时打印的错误，而非引擎的正常产出。整个部署期间，\textbf{没有任何一个自主子任务真正产出过内容}。引擎因此不断重试、生成新的执行命令，形成一个“生成命令—执行失败—再生成”的空循环。

\subsection{小结：一个退化为宏的自主系统}

三个发现指向同一结论：无论是交互式监督路径（退化为反复“继续”），还是自主执行路径（反复失败空转），WhatyTerm 在这次真实部署中\textbf{都未能体现其设计所承诺的“智能编排”}。它在功能上退化为一个昂贵的“自动点继续”宏。值得强调的是，这并非因为系统崩溃或报错——恰恰相反，它“运行良好”地、稳定地做着无用功。这种“看似正常运行、实则空转”的形态，比显式崩溃更危险，也更值得剖析。下一节我们深入代码，解释这一退化为何会稳定发生。
